\begin{frame}{비선형성이 실제로 필요한 순간: $y = \sin(x)$}
  \begin{columns}[T]
    \begin{column}{0.5\textwidth}
      $y = 2x{+}1$은 \textbf{직선} --- 선형 모델 하나로도 풀림.
      \textbf{비선형} $y = \sin(x)$를 같은 2층 MLP로 근사:
      (numpy MLP, 은닉 32유닛, ReLU, SGD lr=0.1, $x \in [-2,2]$,
      512샘플, 800 에포크)
      \begin{itemize}
        \item \textbf{32개의 직선 조각}이 이어붙여 곡선이 됨 ---
          $\sin(x)$ 곡률을 따라 ``계단식 근사''.
        \item 손실: $1.44 \to 0.017 \to 0.008 \to 0.003$
          (1/100/400/800 에포크) --- \textbf{지수적 감소}, 신경망
          학습의 전형적 패턴.
        \item 직선으로 \textit{불가능}했던 근사가 비선형 2층으로
          \textit{가능} --- ``비선형성의 힘''.
      \end{itemize}
      {\footnotesize 주의: 32유닛 계단식 근사는 극값
      ($x = \pm\pi/2$)과 양 끝($x = \pm 2$)에서 오차가 가장 큼 ---
      곡률이 센 곳. 유닛을 늘리면 매끄러워짐.}
    \end{column}
    \begin{column}{0.46\textwidth}
      \centering
      \includegraphics[width=\linewidth]{ch01_2_sin_fitting.png}
      \vskip0.4em
      {\footnotesize 왼쪽: 800에포크 학습 후 근사(파란 원 = 데이터,
      붉은 곡선 = MLP). 오른쪽: 손실 곡선(MSE, 로그 스케일),
      최종 $\approx 0.0031$.}
    \end{column}
  \end{columns}
\end{frame}
